% この本の立場・作り方・参考文献との関係。著作権の方針は docs/knowledge/copyright.md。
\chapter{この本について}\label{ch:preface}

物理を学び直すための\textbf{個人の学習ノート}である。式を追うだけでは想像しにくいところを、
自分たちで動かすシミュレーションと図で確かめながら進む。

\begin{itemize}
  \item \textbf{書き方}: 原稿は生成 AI（Claude）の支援を受けて書いている。主な式と数値は、
    本プロジェクトのシミュレーションとテストで確かめている。それでも誤りが残る可能性がある。
  \item \textbf{参考文献との関係}: 章や節の見出しに、参考にした文献の節番号を示すことがある。
    読み比べるための道しるべであり、本書はその文献の解説書でも副読本でもない。
    著者・出版社とは関係がない。
  \item \textbf{何を写していないか}: 参考文献の文章・図・問題を写して作ったものではない。
    図は本プロジェクトのプログラムで作り、演習は独自に用意した。
  \item \textbf{照合していないこと}: 原典との逐一の照合はしていない。正確な内容は原典で確認してほしい。
\end{itemize}

図・動画の作り方と、各章に対応するコードは、本の中の「シミュレーションで確かめる」に示す。
